\section{Data Structure and Variable Coding}

\subsection{Variable Types}

The survey will contain categorical, binary, ordinal, continuous, and count variables. Variable type will be assigned according to the response format and measurement properties of each survey item.

% TBD: The final variable type assigned to each survey item will be established once the survey
% instrument and response options have been finalised.

\subsection{Coding of Categorical Variables}

Nominal categorical variables will be coded using discrete numeric or labelled values corresponding to the response categories. Categories will remain distinguishable during analysis, and responses such as ``Other,'' ``None,'' and ``Prefer not to say'' will not automatically be treated as equivalent to substantive categories. Free-text responses associated with ``Other'' may be reviewed and recoded where appropriate.

% TBD: The precise coding scheme, category labels, treatment of rare categories, and rules for
% recoding free-text ``Other'' responses will be established before analysis.

\subsection{Coding of Binary Variables}

Binary variables will distinguish the presence or absence of the relevant characteristic or behaviour. Missing, non-applicable, and non-response states will remain distinguishable from substantive responses where required.

For select-all-that-apply questions, each response option will be represented as a separate binary variable indicating whether the respondent selected that option. This will permit multiple characteristics to coexist within a respondent.

% TBD: The final numeric coding convention will be established before analysis.

\subsection{Coding of Ordinal Variables}

Ordinal variables will be coded according to the natural ordering of their response categories while preserving the distinction between adjacent categories. The planned 1--10 preference ratings will retain their numerical ordering, with higher scores representing greater willingness or interest in engaging with the relevant content.

% TBD: The final ordinal response categories, scale anchors, and treatment of ``not applicable,''
% ``prefer not to say,'' and similar responses will be established before analysis.

\subsection{Treatment of Continuous and Count Variables}

Age will be treated as a quantitative variable if collected as an exact value, subject to the final survey design and privacy considerations. Numerical preference ratings will be retained as their observed 1--10 values unless a specific analytical justification for categorisation is established.

Count variables will represent discrete quantities such as the number of fandoms, platforms, Fimfiction groups, conventions, or other activities reported by respondents. These variables will retain their integer values.

% TBD: The final treatment of age, count variables with highly skewed distributions, possible upper or
% lower bounds, and any decisions to categorise continuous variables will be established before analysis.

\subsection{Reference Categories, Where Applicable}

Reference categories will only be required for analyses involving comparisons between categories or statistical models that require a designated reference group. Where subgroup comparisons are conducted, reference categories will be selected based on substantive interpretability, prevalence, or other prespecified methodological considerations rather than being chosen after observing the results.

% TBD: Specific reference categories will be determined if inferential or comparative analyses requiring
% a designated reference category are included.

\subsection{Operational Definitions}

Operational definitions for variables will be based on the survey questions and response categories used to measure the relevant characteristics, behaviours, and preferences. Disability and neurodivergence will remain separate dimensions, and multi-response structures will be retained where respondents may legitimately belong to multiple categories.

% TBD: The final operational definitions and category lists for disability, neurodivergence, minority status,
% active Fimfiction engagement, fandom membership, writer status, and other important variables will be
% established during finalisation of the survey instrument.

\subsection{Derived Variables}

Derived variables may be constructed from raw survey responses where doing so provides a meaningful descriptive measure without obscuring the underlying responses. Potential examples include the number of platforms on which a respondent engages with MLP content, the number of other fandoms reported, the number of Fimfiction groups joined, and the number of disability or neurodivergence categories selected.

Raw responses will be retained wherever possible so that derived variables can be reproduced and alternative classifications can be evaluated.

% TBD: The final set of derived variables, source variables, calculation rules, handling of missing or
% non-applicable responses, and any thresholds used to create derived categories will be specified
% before the final analysis is conducted.
